\section{Conclusion}
\label{sec:conclusion}

We presented Omni-Shield, a local-first multi-modal
PII redaction system that runs entirely on consumer-grade
edge hardware and produces cryptographically verifiable
audit records for every redaction decision.
The system addresses a gap that existing tools leave
open: no prior system simultaneously handles visual
PII (signatures, faces), operates without transmitting
documents to third-party infrastructure, and provides
an independently verifiable audit trail with
zero-knowledge proof of redaction correctness.

\subsection*{Summary of Contributions}

\textbf{Multi-modal nine-agent pipeline.}
We designed and implemented a nine-agent orchestration
pipeline that processes six document formats on a
single 8\,GB GPU under 4-bit NF4 quantisation.
The pipeline achieves 10.42\,s mean latency for image
redaction and under 50\,ms for text and PDF documents.
Ablation confirms that visual detection agents
account for 54\% of total F1 on identity documents,
validating the multi-modal design over text-only
alternatives.

\textbf{Knowledge distillation for document PII.}
We fine-tune Qwen2-VL-2B-Instruct on 2,000
synthetically generated identity documents across six
countries using QLoRA with teacher labels from the
Claude API.
The primary contribution of fine-tuning is output
format consistency: the adapter trains the model to
emit structured \texttt{LABEL\,|\,VALUE\,|\,BBOX}
triples reliably, enabling downstream pipeline agents
to parse and act on VLM output without heuristic
recovery.
Face detection generalises as an emergent capability
post fine-tuning, despite faces not being explicitly
labelled in the teacher annotations.

\textbf{Cryptographic audit trail with ZK-SNARK.}
Every redaction is anchored on a local Ethereum
blockchain using a salted SHA-256 commitment that
prevents linkability attacks.
A Groth16/BN128 proof generated by a ZoKrates circuit
attests that the committed PII values are consistent
with the document hash, without revealing what was
redacted.
Blockchain verification completes in 1.2\,ms.
ZK proof verification completes in 22.1\,ms.
Proof generation takes 20.3\,s but is fully
asynchronous, adding no user-facing latency.
The system is the first, to our knowledge, to apply
ZK-SNARKs to document redaction provenance.

\textbf{Comprehensive evaluation on synthetic and
real documents.}
We evaluate against Microsoft Presidio and Llama~3~8B
on 150 annotated synthetic cards and 19 real specimen
documents.
Omni-Shield is the only system to detect signatures
(F1\,=\,0.870) and faces (F1\,=\,0.905 on real
documents), and achieves Macro F1\,=\,0.269 on real
specimen cards --- higher than the synthetic test
result, indicating positive transfer from the
synthetic training distribution.

\subsection*{Key Findings}

Three findings from this work have broader implications
for the design of privacy-preserving document systems.

\emph{Text-only systems are insufficient for identity
document redaction.}
The strongest text-only baseline (Llama~3~8B,
Macro F1\,=\,0.270) scores zero on signatures and
faces --- the two most legally significant visual PII
categories.
Any deployment that relies on text extraction alone
will systematically fail to redact the most visually
distinctive personal information on identity documents.

\emph{Local-first deployment is not a performance
concession.}
Despite operating on a consumer GPU at 4-bit
quantisation, Omni-Shield outperforms a mature
cloud-backed NLP system (Presidio, Macro F1\,=\,0.057)
by a factor of 3.3 on the same test set, and
achieves competitive latency (10.42\,s per image)
against cloud vision APIs that require network
round-trips for every document.
The privacy premium of keeping documents on-device
amounts to 0.089\,kWh per 1,000 documents
(approximately \rupee0.72 at Indian electricity
rates) --- a quantifiable and modest overhead.

\emph{Multi-modal agent pipelines are more robust
than monolithic models.}
The CriticAgent pattern, ablation over agent
configurations, and the TEXT-FIRST bounding box
refinement strategy all demonstrate that decomposing
the redaction task into specialised agents --- each
with a defined input, output, and failure mode ---
produces a system that is more debuggable,
more ablatable, and more improvable than a
single model would be.
Each agent's contribution is measurable; underperforming
agents can be replaced without redesigning the pipeline.

\subsection*{Limitations}

The primary limitations of the current system are:
low recall on ID numbers (F1\,=\,0.096), a synthetic-
only training distribution with documented transfer
gaps for diverse name typography, a trusted setup
requirement for the ZK circuit that weakens the
proof's soundness guarantee in the absence of a
multi-party ceremony, and Ethereum mainnet gas costs
that require L2 deployment for economically viable
production use.
These are discussed in detail in
Section~\ref{sec:limitations}.

\subsection*{Future Work}

The most impactful near-term improvements are:
(i)~replacing the ZK trusted setup with a
public MPC ceremony or transparent proof system;
(ii)~expanding the training corpus with 500
additional cards targeting ID number format diversity
to address the recall gap;
(iii)~combining Omni-Shield's visual detection with
a text-specialised model (such as Llama~3) in a
hybrid pipeline that achieves high F1 across both
visual and text PII categories;
and (iv)~deploying the blockchain contract on
Polygon L2 for production-ready gas economics.

\subsection*{Reproducibility}

The fine-tuned QLoRA adapter is available at
\url{https://huggingface.co/[anonymous]/omni-shield-qwen2vl-pii-adapter}.
The evaluation harness, annotation pipeline,
synthetic dataset generator, and all experiment
scripts are available at
\url{https://github.com/[anonymous]/omni-shield}.
A Docker container with GPU support is provided
for one-command environment setup.
All reported numbers are reproducible from the
published code and data using the commands documented
in the repository README.

\subsection*{Ethical Considerations}

Omni-Shield is designed to protect privacy, not
to circumvent it.
The system does not retain documents after
processing, does not transmit data off-device,
and does not log PII values in plaintext beyond
the masked audit record.
The synthetic training data contains no real
personal information.
Specimen documents used in real-document evaluation
are officially published educational materials
containing no real PII.
The blockchain audit trail is designed to provide
accountability for \emph{operators} of redaction
systems, not to surveil the individuals whose
documents are being processed.
We release all artifacts under Apache~2.0 licence.
